\section{Spinors}

Why should a theory of feeling care about spinors at all? Because spinors are what matter is. Every particle that builds an atom, the electron and the quarks inside the proton, is a spinor, a thing of half-integer spin, and it is their spinor nature that gives matter its solidity, the exclusion that forbids two electrons the same state and so keeps the world from collapsing to a point. Light and the forces are not spinors, they are vectors, and a universe able to hold only vectors would be radiation with nothing to radiate from, no matter, no chemistry, no observers. So a substrate that means to grow a real world has one hard requirement ahead of all the others. It must be able to carry a spinor. Most discrete spaces cannot, which is exactly why building matter on a lattice has been so hard. The remarkable thing about this one is that it can, and not by adding spin as an ingredient but by having it forced on it by the shape of the dock.

The strangest fact about matter is that an electron must be turned around twice to come back to itself. One full turn leaves it changed, its sign flipped, and only a second turn restores it. This is spin one-half, normally a fact about continuous rotations and the double cover of the rotation group by $\mathrm{SU}(2)$. The substrate carries it with no continuum at all, and the way it does is worth showing in full, because it is the cleanest case of physics being forced by the dock rather than fitted to it.

The method starts by noticing that the dock's twenty-four directions are not merely a set of arrows, they are a group. Written as unit quaternions they close under multiplication, the binary tetrahedral group of twenty-four elements, and written as coordinate vectors they are the same twenty-four $D_4$ roots, the points with two entries $\pm 1$ and two zero. One object with two faces, a group of rotations and a root system at once.

That group is exactly the double cover. The unit quaternions are $\mathrm{SU}(2)$, which wraps twice around the rotation group, so each ordinary rotation lifts to two quaternions, one the negative of the other. The two lifts act differently on different things, and that difference is the whole of spin. A vector is turned by conjugation, sandwiched between a quaternion and its inverse, where the two lifts cancel and a single full turn brings the vector home. A spinor is turned by plain multiplication, the half-angle action, where the two lifts do not cancel. Take a quaternion whose square is $-1$, a quarter turn's worth, and apply it twice to make one full turn of a spinor. The result is not the spinor but its negative. Apply it four times, two full turns, and the spinor returns. The minus sign at $2\pi$ and the return at $4\pi$ are not argued from a picture, they are read straight out of the twenty-four-element multiplication table by composing its elements, exact and continuum-free.

The same group hands over the rest. Inside the twenty-four sit eight that form the quaternion group, the vector octet, and multiplying them by a primitive cube root of unity, the quaternion $(-1 + i + j + k)/2$, cycles to a second octet and then a third. Three octets of eight, $8v$ and $8s$ and $8c$, a vector and two spinors of opposite hand, and the threefold cycle that relates them is the triality met in the geometry, the seed of the three generations. Build the Dirac and Pauli matrices on top and they inherit the same sign, the Clifford relations holding exactly and the chirality operator anticommuting with the mass term, so the full machinery of the relativistic electron stands on the one finite group.

\begin{result}[Spinors on the dock]
The dock's twenty-four directions form the binary tetrahedral group, the double cover of the rotation group, so a $2\pi$ rotation acts as $-1$ and a $4\pi$ rotation as $+1$, exactly, by composition in a finite group. The directions split by triality into $8v + 8s + 8c$, a vector and two spinors. \cite{pollard2026vibetest}
\end{result}

The contrast with the rejected substrate is sharpest here, and it is a fact of representation theory rather than taste. The dock of $\{5,3,4\}$ is twelve directions, the vertices of an icosahedron, and their rotation symmetry is the icosahedral group. Its action on those twelve directions breaks into pieces of integer spin only, a singlet, two triplets, and a quintet, $1 + 3 + 3' + 5$, with no spinor among them. A spinor exists for that group, but only in its own double cover, the icosian group, which does not sit on the twelve directions the way the binary tetrahedral group sits on the twenty-four. On $\{5,3,4\}$ the spinor is absent from the directions a tone actually moves along. On $\{3,4,3,4\}$ it is built into them. That single difference, a spinor present in the dock or missing from it, is why matter can live on the one substrate and not the other.
